% Scientific description for the annual policy-association results.
% The coefficient and weather tables report complementary terms.
\subsection{Annual policy associations}
\label{sec:si-annual-policy-associations}

The annual comparison estimates conditional associations between documented
policy functions and three source-defined service measures.  It is distinct
from the quarterly facility--activity model and is not a policy mediation
analysis.  Let $Q$ denote compatible-week recorded charging activity, $E$ the
compatible-week cumulative facility exposure, and $W$ the number of compatible
weeks.  The outcomes are $A=Q/W$, $N=E/W$ and $I=Q/E$, with fitted outcomes
$\log A$, $\log N$ and $\log I$.  Thus $N$ is not a new-build count, rated
power or connector capacity, and the annual comparison does not replace source
denominators with calendar-year exposure.

\subsection{Policy support and indicators}
The association support is defined jointly by source observations and policy
coverage.  The annual city--year registry contains 334 cities.  An association
row requires a numerically nonmissing source facility-additions measure with
one of six observed statuses (observed complete, observed partial, observed
zero, observed, complete or partial), observed compatible $Q$, positive
compatible $I$, and finite values for all four annual weather controls.
Prior-year construction-support and operation-policy counts must be known under
both effective-year assignment and the effective/publication-year rule in
Section~\ref{sec:si-policy}.
These completeness conditions define common support across the two timing conventions;
the reported construction-support regressor $B$ uses only the latter rule.  Rows in this support have positive
compatible-week $Q$, $E$ and $W$; $R/S$ coverage is observed locally and its
timing status is known.  A documented absence of a qualifying requirement
within observed local coverage is an observed zero.  Unknown city coverage or
timing is not recoded as zero.

The facility-observation condition is a source-status condition, not a
construction quantity condition.  It accepts a row only when the source's
facility-additions measure is numeric and nonmissing and its status is one of
the six observed statuses listed above.  The resulting support contains 269
observed-complete rows and 56 observed-zero rows.  Missing or uncovered rows,
pre-entry structural rows, source-coded unknowns and other unlisted statuses
are excluded.  Thus the condition certifies that the source facility-additions
measure has an interpretable observation status; it does not require $E>0$ and
does not define the facility exposure used in the outcomes.  The outcome
definitions remain $E$ equal to the sum of the cumulative facility indicator
over compatible weeks and $W$ equal to the number of compatible weeks, so that
$N=E/W$ is a mean facility indicator rather than a period-additions count.
The common eligible pool additionally requires $Q>0$, $E>0$ and $W>0$,
followed by the stated year and complete COVID-exposure conditions.  These
conditions yield the 325-row, 94-city policy support.

The comparison is restricted to 2016--2022 and rows with complete source-date
COVID exposure.  Pre-pandemic structural zeros in that exposure series remain
complete observations; source gaps are not filled.  The association support contains
325 city--year rows, 94 cities and 25 province clusters.  Seventy-five cities
contribute at least two years and 19 are singletons; 10 cities switch in the
information/supervision indicator and 6 switch in the operation/maintenance
indicator.  The COVID exposure consists of 214 source-measured complete rows and 111 pre-pandemic
structural-zero rows.  These counts define the estimand and do not imply a
balanced panel.

For city $i$ and calendar year $t$, the indicators are measured at the end of
the preceding year.  $R^{\mathrm{obs}}_{i,t-1}$ records an observed
information or supervision function, including platform connection, data
reporting, information interoperability, state monitoring, and data collection
or sharing.  $S^{\mathrm{obs}}_{i,t-1}$ records an observed operation or
maintenance responsibility, including maintenance operations, fault-response
repair, availability or normal operation, emergency continuity and related
operator qualifications.  These indicators describe documented function and
timing; they are not continuous policy-strength scores and do not measure
implementation intensity, compliance, service volume or realized delivery.

\subsection{Specifications and inference}
For $Y\in\{A,N,I\}$, the two annual specifications are
\begin{equation}
\log(Y_{it}) = \alpha_i+\tau_t
 + b_R R^{\mathrm{obs}}_{i,t-1}
 + b_S S^{\mathrm{obs}}_{i,t-1}
 + b_B B_{i,t-1}
 + \gamma_T T_{it}+\gamma_P P_{it}
 +\gamma_V V_{it}+\gamma_H H_{it}
 + d_C\,C_{it}/10 + \varepsilon_{it},
\label{eq:annual-policy-association}
\end{equation}
where $B$ is the construction-support document count assigned by the annual
effective/publication-year rule in Section~\ref{sec:si-policy}, and $B_{i,t-1}$ counts documents assigned
to the preceding calendar year; it is not a cumulative stock.  The annual $B$
assignment is separate from the quarterly date-interval classification.  $T$, $P$, $V$ and $H$ are annual temperature,
precipitation, wind speed and
relative humidity, and $C$ is the current-year source-date COVID stringency
mean.  The base specification omits the final COVID term; the adjusted
specification includes it.  These four weather variables are annual policy
controls, not the five-variable facility elasticity set or quarterly weather
values.  City and calendar-year fixed effects are included in every outcome
model, with equal weight on log-transformed city--year observations.

Inference uses province-clustered standard errors with the finite CR1
correction
\[
 \frac{G_p}{G_p-1}\frac{n-1}{n-K},
\]
where $G_p=25$, $n=325$, and $K$ is the full design rank including fixed
effects and displayed regressors.  The reference distribution is
$t(G_p-1)=t(24)$ and intervals are two-sided 95\% intervals.  The base and
COVID-adjusted designs have full ranks 107 and 108 on the same 325-row
support.  The policy-association table reports every policy term, lagged
construction term and response-index term with its estimate and interval.
Table~\ref{tab:si_policy_weather} reports all four weather adjustment
coefficients for both specifications on their raw annual source scales:
temperature in degrees Celsius, precipitation in mm/year, wind speed in m/s
and relative humidity in percent.  The full covariance blocks among substantive
terms are retained, so linear contrasts use covariance rather than rounded cell
values.

\subsection{Complete policy results and interpretation}
The main policy table presents the core R/S associations (Table~\mainref{tab:policy_associations});
the complete R/S results are retained in Table~\ref{tab:si_policy_association}
for $\ln A$, $\ln N$ and $\ln I$ under both specifications, together with
the adjustment terms in the companion weather table.  The information/
supervision association with $\ln A$ is positive in both specifications
(1.770, 95\% CI 0.413 to 3.127; and 1.739, 95\% CI 0.383 to
3.096).  The corresponding intervals for $\ln N$ and $\ln I$ include zero,
so this comparison does not identify a common shift across the three service
measures.

All operation/maintenance confidence intervals include zero; their widths vary
across outcomes.  An interval spanning zero does not establish that a coefficient equals zero.

The COVID-adjusted response-index terms are adjustment terms in the stated
specification, not additional policy contrasts; the lagged construction count
and weather terms have the same status.  The information/supervision--activity
association is compatible with a difference in recorded activity conditional on
the stated support and controls.  It can also reflect reporting, monitoring or
information visibility, so it is not evidence by itself of additional delivered
service or a causal policy effect.  The documented functions and the quarterly
governance-support comparison have separate evidential roles; no coefficient is
interpreted as a net policy effect or joined to the governance comparison as a
mediation chain.
